% ============================================================
%  最小编译自检模板 —— 用于验证 xelatex + ctex + 数学环境可用
%  compile_pdf.py 会自动用真实内容替换，此文件仅供 smoke test
% ============================================================
\documentclass[11pt]{ctexart}
\usepackage{amsmath,amssymb,amsthm}
\usepackage{geometry}
\geometry{a4paper, margin=2.5cm}
\usepackage{enumitem}
\usepackage{hyperref}
\usepackage{tcolorbox}
\tcbuselibrary{theorems, breakable, skins}

% --- 例题/概念/公式盒子的统一定义（compile_pdf.py 生成的文档会复用这套）---
\newtcbtheorem[number within=section]{examplebox}{例题}{
  colback=blue!4, colframe=blue!45!black, fonttitle=\bfseries,
  breakable
}{ex}
\newtcbtheorem[number within=section]{conceptbox}{概念}{
  colback=green!4, colframe=green!40!black, fonttitle=\bfseries,
  breakable
}{cpt}

\title{编译自检 \\ Compilation Smoke Test}
\author{lecture-organizer}
\date{}

\begin{document}
\maketitle

\section{中文与公式}
本章覆盖\textbf{数字通信}（Digital Communications）中的关键概念，
例如信噪比 SNR $= \frac{E_b}{N_0}$，以及误码率。

\begin{equation}
\text{BER} = Q\left(\sqrt{\frac{2E_b}{N_0}}\right)
\end{equation}

\begin{itemize}
  \item 复数信号：$z = a + jb$，模 $|z| = \sqrt{a^2 + b^2}$
  \item 傅里叶变换：$\displaystyle X(f) = \int_{-\infty}^{\infty} x(t)\,e^{-j2\pi ft}\,dt$
\end{itemize}

\section{例题}
\begin{examplebox}{已知余弦信号的频谱}{cos-spectrum}
已知 $x(t) = \cos(2\pi f_0 t)$，求其频谱 $X(f)$。
\end{examplebox}

\end{document}
